\section{Validation}\label{validation}

\subsection{Automatic Testing}\label{automatic-testing}

Testing is organised in three tiers:

\begin{sloppypar}
\paragraph{Unit tests.} Every one of the nine backend components (the gateway
and the eight domain services) has its own, fully isolated Jest
configuration and a \texttt{\_\_tests\_\_} folder, for a total of \textbf{44
test files}. Consistent with the hexagonal layering of
\cref{technological-details}, these tests never touch a real MongoDB or
Redis: each service provides a \texttt{\_\_tests\_\_/helpers/fakes.ts} with
in-memory implementations of its domain ports so a use-case test real business logic
against a deterministic double instead of infrastructure. Two kinds of
targets are covered: \emph{use-cases} and \emph{pure domain functions}.
HTTP adapters are also unit-tested, with \texttt{supertest} driving an
Express app built purely from fakes (\texttt{buildTestApp.ts}), so a router
test verifies status codes, validation and role/ownership rejection without
any real dependency. Coverage is measured for an informative purpose, since there is no CI.
\end{sloppypar}

\paragraph{Domain-level concurrency tests.} The most distributed-systems-relevant
unit tests are cadService's \texttt{collabConcurrency.test.ts} and
\texttt{collabCheckpointRecovery.test.ts}, which exercise
\texttt{InMemoryCollabSessionService} (\cref{fault-tolerance}) directly
against fakes. The first covers twelve scenarios of the Lamport/LWW conflict
rule: a higher Lamport timestamp wins, a tie is broken by actor id,
independent fields converge independently, a duplicate \texttt{opId} is
acknowledged without being re-applied, an old \texttt{baseVersion} is
rejected, and the outcome is independent of the arrival order of two
concurrent operations.
The second test drives a full crash/recovery cycle: it checkpoints a
session, creates a \emph{brand-new} \texttt{InMemoryCollabSessionService}
instance sharing the same checkpoint/event-log fakes (simulating a process
restart), calls \texttt{recoverSessions()}, and asserts the rebuilt state
(name, category, version, participants) is identical to the pre-crash state.

\begin{sloppypar}
\paragraph{Integration tests (``e2e'').} A separate, top-level \texttt{e2e/}
package largely re-exercises the same business rules already covered by unit
tests, but as black-box HTTP calls against the \emph{real}, dockerised stack
(real MongoDB/Redis instead of fakes).
Its \texttt{globalSetup.js} polls the gateway's aggregated
\texttt{/health} until the stack is up, obtains a
guest JWT, and idempotently seeds an \texttt{ADMIN} dev user before the suite
runs. Eight spec files (\texttt{auth}, \texttt{cad}, \texttt{cart},
\texttt{catalog}, \texttt{chatbot}, \texttt{notifications}, \texttt{order},
\texttt{reporting.integration.test.ts}), each tagged \texttt{[INT]}, cover
registration/login/session/logout, configurator, cart/catalog/order CRUD,
chatbot Q\&A and the admin-only reporting endpoint, with ownership isolation
(403/404 on another user's resource) and role rejection asserted throughout
(directly validating NFR5). Two things go beyond a mere replay with real
infrastructure, since they involve behaviour unit tests cannot see by
construction: the full configurator flow for both the TONDO and the INTELLIGENTE product family 
(environment $\to$ category $\to$ column
plan $\to$ design $\to$ finalize, asserting the returned BOM) and catalog's
\texttt{[DS]} case (\texttt{PUT /catalog/:id} against an old version $\to$
\texttt{409 VERSION\_CONFLICT}) both exercise real persistence end to end.
\end{sloppypar}

\begin{sloppypar}
\paragraph{How to run them.} Every service exposes \texttt{npm test} and \texttt{npm run test:coverage}; 
the workspace
root aggregates all nine with \texttt{npm run test:backend}, the black-box
suite with \texttt{npm run test:e2e} (which requires
\texttt{docker compose -f docker-compose.dev.yml up} to already be running),
and \texttt{npm run test:baseline} chains both. There is currently
\textbf{no continuous-integration pipeline} running any of this
automatically on push or before merge.
\end{sloppypar}

\subsection{Acceptance test}\label{acceptance-test}

Two kinds of manual verification complement the automated suites, since
neither is easily expressed as a Jest assertion.

\paragraph{Manual UI testing.} The Vue frontend has no automated test suite
(\cref{technological-details}), so every user-facing flow described in
\cref{user-guide} as the registration/guest access, the full configurator
wizard for each of the four product families, the collaborative session
(joining from a second browser tab/profile to see real-time synchronisation),
cart/checkout, the notification bell and toast pushes, the chatbot, and the admin screens,
was exercised manually against the Docker Compose stack.

\paragraph{Production-like deployment as a test.} Following the hint that
deployment automation doubles as a test of a production-like environment,
the two distributed-systems properties that automated tests can only
approximate with fakes were also demonstrated manually on the Kubernetes
deployment (\cref{deployment}): killing the current CAD MongoDB primary
(\texttt{kubectl delete pod mongo-cad-0}) and observing the replica set
elect a new primary and writes resume; and killing the running
\texttt{cad-service} pod (\texttt{kubectl delete pod -l app=cad-service}) and
observing an open collaborative session survive via checkpoint/replay,
validating NFR4 against a real process crash and
Kubernetes-driven restart, not just the
\texttt{collabCheckpointRecovery} unit test.
